\documentclass[11pt,a4paper]{article}

% ---------- packages ----------
\usepackage[T1]{fontenc}
\usepackage[utf8]{inputenc}
\usepackage{amsmath}
\usepackage[margin=2.3cm]{geometry}
\usepackage[table]{xcolor}
\usepackage{booktabs}
\usepackage{array}
\usepackage{tabularx}
\usepackage{enumitem}
\usepackage{titlesec}
\usepackage{fancyhdr}
\usepackage{graphicx}
\usepackage[hidelinks]{hyperref}

% ---------- colours ----------
\definecolor{navy}{HTML}{000000}
\definecolor{accent}{HTML}{000000}
\definecolor{lightblue}{HTML}{F2F2F2}
\definecolor{lightgreen}{HTML}{E8E8E8}
\definecolor{amber}{HTML}{ECECEC}
\definecolor{greyln}{HTML}{808080}

% ---------- section styling ----------
\titleformat{\section}{\Large\bfseries\color{navy}}{\thesection.}{0.6em}{}[{\color{accent}\titlerule[1pt]}]
\titleformat{\subsection}{\large\bfseries\color{black!80}}{\thesubsection}{0.6em}{}
\titlespacing*{\section}{0pt}{14pt}{6pt}

% ---------- header / footer ----------
\pagestyle{fancy}\fancyhf{}
\renewcommand{\headrulewidth}{0.3pt}
\lhead{\small\color{navy}AMR-X Prototype V1}
\rhead{\small\color{navy}Payload-Deck Suspension Study}
\cfoot{\small\color{black!55}Page \thepage}

\newcolumntype{Y}{>{\raggedright\arraybackslash}X}
\newcommand{\prio}[1]{\colorbox{amber}{\parbox{\dimexpr\linewidth-2\fboxsep}{#1}}}

\setlist[itemize]{leftmargin=1.3em,itemsep=2pt,topsep=3pt}

% ---------- title ----------
\begin{document}
\begin{titlepage}
\thispagestyle{empty}
\vspace*{2.2cm}
\noindent\rule{\linewidth}{1.2pt}\\[10pt]
{\small\bfseries\color{black} AMR-X PROTOTYPE V1 \quad$\cdot$\quad MECHANICAL ENGINEERING}\\[14pt]
{\color{black}\Huge\bfseries Payload-Deck Suspension}\\[8pt]
{\color{black}\Large Shock Isolation for the Payload Plate: Options, Trade-offs, Recommendation and Procurement}\\[12pt]
\noindent\rule{\linewidth}{0.5pt}\\[10pt]
{\normalsize 100\,kg payload \quad$\cdot$\quad 600--800\,mm deck plate \quad$\cdot$\quad impact-dominated duty \quad$\cdot$\quad vibration + curb-drop secondary}
\vfill
\noindent\rule{\linewidth}{0.5pt}\\[4pt]
{\footnotesize Engineering design report \hfill Revision 1}
\end{titlepage}

% =====================================================================
\section*{List of Abbreviations}
\begin{center}\small
\renewcommand{\arraystretch}{1.25}
\begin{tabularx}{\linewidth}{@{}p{2.3cm}Y@{}}
\toprule
\textbf{Abbreviation} & \textbf{Meaning} \\
\midrule
AMR & Autonomous Mobile Robot \\
AMR-X & Designation of this prototype platform (the robot studied in this report) \\
AGV & Automated Guided Vehicle \\
\rowcolor{lightblue} WRI & Wire Rope Isolator (looped stranded-cable shock/vibration mount) \\
PU & Polyurethane (elastomer / wheel material) \\
\rowcolor{lightblue} CoG & Centre of Gravity \\
RFQ & Request for Quotation \\
\rowcolor{lightblue} MOQ & Minimum Order Quantity \\
NVH & Noise, Vibration and Harshness \\
\rowcolor{lightblue} USPTO & United States Patent and Trademark Office \\
MIL-STD & United States Military Standard (e.g.\ MIL-STD-810 environmental) \\
\midrule
\multicolumn{2}{@{}l}{\textit{Units}} \\
N / kN & Newton / kilonewton (force) \\
\rowcolor{lightblue} J & Joule (energy) \\
kg & Kilogram \\
\rowcolor{lightblue} mm & Millimetre \\
Hz & Hertz (natural frequency) \\
\rowcolor{lightblue} g & Acceleration relative to gravity ($\approx$9.81\,m/s$^2$) \\
m/s & Metres per second \\
\bottomrule
\end{tabularx}
\end{center}

% =====================================================================
\section{Why this question exists}
The payload deck currently bolts \textbf{rigidly} to the lift/frame. That is fine until the moment a
$\approx$100\,kg payload is \emph{set down} on it. The kinetic energy of that placement has to go
somewhere, and on a rigid deck it goes straight into the frame, the fasteners and the lift mechanism as
a sharp force spike.

The governing relation is simple and unforgiving:
\[
\textbf{peak force into the frame} \;\approx\; \frac{\text{impact energy}}{\text{stroke}}.
\]
A feel for the numbers: 100\,kg arriving with even $\approx$25\,mm of drop is $\approx$25\,J. Absorbed
over 25\,mm of suspension travel that averages $\approx$1\,kN ($\approx$100\,kgf). Absorbed over the
1--2\,mm of ``give'' a rigid bolted deck has, the same 25\,J becomes \textbf{12--25\,kN} slammed into the
structure. Interposing controlled travel buys an order-of-magnitude ($\approx$10$\times$) reduction in
peak force. \textbf{That force reduction --- not ride comfort --- is the whole point of this mechanism.}

\textbf{The duty, in priority order.}
\begin{itemize}
\item \textbf{Primary: placement shock.} Protect the robot when the payload lands on the deck.
\item \textbf{Secondary: floor vibration and curb drops} transmitted up through the chassis while driving.
\end{itemize}

\textbf{The catch a spring alone cannot solve.} A pure spring (rubber, coil or air) \emph{stores} the
placement energy and throws it straight back --- the payload bounces and you get repeated impacts. For
placement shock the mechanism must \textbf{dissipate} energy, not just deflect. This single fact drives
the entire ranking that follows.

\textbf{Design target.} Impact-dominated, with the payload varying $0\rightarrow100$\,kg every cycle, on
an indoor warehouse floor (small repeated unevenness --- joints, ramps, dock-plate lips --- not rough
terrain). Crucially the deck must still present a \emph{repeatable height} for docking and pick/place;
production AMRs in this class hold $\approx\pm5$\,mm positioning, and AGV practice warns that an
unconstrained compliant deck ``floats'' its height and breaks dock/sensor alignment~\cite{abb,casterhq}.
The real question is therefore: \textbf{how do we put controlled, damped travel between deck and frame
without losing deck-height repeatability.}

% =====================================================================
\section{Candidate mechanisms}
Six approaches span the spectrum from doing nothing to a fully guided, damped platform.

\begin{center}\small
\renewcommand{\arraystretch}{1.25}
\begin{tabularx}{\linewidth}{@{}p{3.5cm}Y c c c@{}}
\toprule
\textbf{Option} & \textbf{How it works} & \textbf{Energy handling} & \textbf{Complexity} & \textbf{Fit} \\
\midrule
A. Rigid deck & plate bolted hard to frame & Poor (no stroke) & Lowest & Gentle hand-loading only \\
B. Elastomeric mounts & rubber sandwich/cone pads & Fair; rebounds & Low & Vibration, light shock \\
C. Coil-spring isolators & steel springs in housings & Fair; stores \& rebounds & Low & Low-freq vibration \\
D. Air springs / pneumatic & rubber bellows on air & Good vibration; needs air & Med & Vibration isolation \\
\rowcolor{lightgreen} E. Wire rope isolators & looped cable, high internal (hysteretic) damping & Good; damps, low rebound & Low & \textbf{Our case} \\
\rowcolor{lightgreen} F. Guided deck + hydraulic damper & plate on linear guides + shock absorbers & Excellent; near-zero rebound & Med--High & \textbf{Our case (worst drops)} \\
\bottomrule
\end{tabularx}
\end{center}

\noindent\textbf{The two that fit us are E and F.}
\textbf{E} (wire rope isolators) puts high energy absorption per millimetre of travel into a single
bolt-on metal part: the cable strands rub and dissipate the impact instead of springing it back, and the
same part also covers the secondary vibration and curb-drop duty in all three axes. \textbf{F} (a plate
on linear guides with hydraulic shock absorbers) gives the cleanest possible \emph{no-rebound} stop for
genuine drops, because a hydraulic damper converts the impact almost entirely to heat at a near-constant
low force. In practice the two combine well --- wire rope isolators as the primary mounts, with a rubber
end-stop, or the guided-deck variant when loads are dropped rather than placed.

Options B, C and D are kept as the \emph{vibration} layer, not the impact layer: a single-rate rubber,
coil or air spring sized to catch a 100\,kg drop is far too stiff when the deck is empty, and rubber can
take a set or tear under repeated heavy impacts.

% =====================================================================
\section{Real-product precedents}
\textbf{Option E --- wire rope isolators.}
These are the defence- and aerospace-standard answer to \emph{repeated} shock precisely because they
damp in all three axes, are unaffected by oil, temperature and ageing, and are maintenance-free; the
standard catalogue lines are qualified to environmental standards such as MIL-STD-810~\cite{enidine}.
The recurring design rule is to mount them at $\approx 45^{\circ}$ to the load, which equalises the
stiffness across all three axes so one part handles the vertical placement shock \emph{and} the
horizontal curb/stop loads~\cite{uspto45a,uspto45b}. The category is already used in automated-warehouse
equipment~\cite{enidine}.

\textbf{Option F --- guided deck with hydraulic shock absorbers.}
Industrial self-compensating shock absorbers convert impact to heat and do not rebound. The decisive
feature for us is that the payload energy changes every cycle as the load varies $0\rightarrow100$\,kg,
and self-compensating units ``react to changing energy conditions without adjustment'', covering
effective weights from $\approx$1\,kg to 1{,}550\,kg with a side-load adapter for impact angles up to
$25^{\circ}$~\cite{acesc,aceheavy}. Mounted under a linearly-guided plate with light return springs,
they form the no-bounce drop-catcher of Solution~F.

\textbf{Why not a bare spring.}
Coil, rubber and air mounts are excellent \emph{vibration} isolators --- air springs in particular reach
very low natural frequencies~\cite{firestone,mason} --- but used alone for placement shock they store the
energy and bounce the load. And on an AGV/AMR specifically, soft unconstrained compliance under the deck
varies the deck height and ``breaks sensor and dock-plate alignment'', which is why such compliance is
normally reserved for manual carts unless its travel is constrained~\cite{casterhq}.

% =====================================================================
\section{Recommendation: Solution E (wire rope isolators)}
\textbf{We recommend Solution~E} --- mounting the deck plate on wire rope isolators at the four corners.
It attacks our exact failure mode (placement energy spiking into the frame) at the source, it damps
rather than rebounds, and it delivers the primary \emph{and} both secondary duties in one
maintenance-free, bolt-on metal part with no air supply or power. It is also the only one of the two
fitting options available as a \emph{finished, purchasable component}, which de-risks and accelerates the
prototype.

\subsection{The trade-off you must accept}
\prio{\textbf{A compliant deck trades structural protection against deck-height repeatability.}
The same travel that protects the frame also lets the deck move under load, so dock/pick height can drift
and the payload can sway in turns. Accept this by \textbf{constraining the travel}: hard limit stops plus
a rubber end-stop, modest stroke ($\sim$15--25\,mm), and isolators placed symmetrically about the payload
CoG so each corner carries an equal $\approx$25\,kg share. If a \emph{precise} docking height is required,
decouple the two needs in time --- soft-isolate while driving, then mechanically lock or clamp the deck
during pick/place and docking.}

\medskip
\noindent If loads can be genuinely \emph{dropped} (not gently lowered), step up to Solution~F: a guided
deck on linear bearings with hydraulic self-compensating absorbers and light return springs. It gives the
cleanest no-rebound stop, but --- like a fabricated rocker in the drivetrain study --- it is an
\emph{assembly}, not a single purchasable kit.

% =====================================================================
\section{Specifications to preserve or improve}
Any deck-suspension solution must \textbf{meet or beat} the rigid baseline on the points below. Targets
derived from the impact and packaging study:

\begin{center}\small
\renewcommand{\arraystretch}{1.25}
\begin{tabularx}{\linewidth}{@{}p{5.0cm}p{4.0cm}Y@{}}
\toprule
\textbf{Parameter} & \textbf{Current (rigid deck)} & \textbf{Requirement} \\
\midrule
Static payload & 100\,kg & support 100\,kg + dynamic margin \\
Per-corner static load & $\approx$25\,kg (4 mounts) & isolator at $\sim$50\% of rating (shock headroom) \\
Vertical travel & $\approx$0 (rigid) & $\sim$15--25\,mm \emph{controlled}, with hard stops \\
Peak force into frame & 12--25\,kN on a drop & cut $\approx$10$\times$ by adding stroke \\
Rebound after placement & n/a & damped, no bounce (dissipative) \\
Load variability & n/a & cope $0\rightarrow100$\,kg per cycle \\
Deck-height repeatability & fixed by bolt-up & maintain for docking (limit stops / lock) \\
Environment & indoor warehouse & oil/temp/age-stable, maintenance-free \\
\bottomrule
\end{tabularx}
\end{center}

\noindent A 4-corner wire rope set clears these with margin; the items needing active confirmation in the
RFQ are the \textbf{drop height / placement velocity} (which fixes the energy and therefore the isolator
size and count) and the \textbf{travel-limit / end-stop} scheme that protects deck-height repeatability.

% =====================================================================
\section{Ready-to-buy products suited to our use case}
The following are appropriate for a 100\,kg payload deck (with the noted configuration). All are
catalogue parts with low MOQ.

\begin{center}\small
\renewcommand{\arraystretch}{1.3}
\begin{tabularx}{\linewidth}{@{}p{3.2cm}Y p{2.5cm}@{}}
\toprule
\textbf{Product} & \textbf{Why it fits us / what to configure} & \textbf{Type} \\
\midrule
\rowcolor{lightblue}
Enidine WR6-700-10~\cite{wr6} &
$\approx$50.8\,kg/unit, $\approx$43\,mm tall. Four corners at $45^{\circ}$; at 25\,kg/corner it sits at
$\sim$50\% rating, leaving compression travel for the impact. \textbf{Primary pick.} & Wire rope isolator \\
Enidine WR8-400-08~\cite{wr8} &
$\approx$68.9\,kg/unit, $\approx$54\,mm tall. Step up to this (or run six WR6) if placement drops are
harsh. & Wire rope isolator \\
\rowcolor{lightblue}
ACE SC$^2$ self-compensating absorbers~\cite{acesc} &
Effective weight 1--1{,}550\,kg, side-load adapter to $25^{\circ}$, self-compensates as load varies
$0\rightarrow100$\,kg. The energy dump for Solution~F. & Hydraulic damper \\
Parker LORD Flex-Bolt sandwich (small group)~\cite{lord} &
$\approx$3--66\,kg/mount; cheap, simple vibration layer to add beneath/alongside the primary mounts. &
Elastomeric mount \\
ACE PAL-3~\cite{acepal} &
Auto-levelling air spring, $\leq$385\,kg/spring; adds constant deck level + isolation if levelling is
needed (requires air). & Pneumatic levelling \\
\rowcolor{lightblue}
PGFUN 10-loop WRI~\cite{pgfun} &
Low-cost 10-loop stainless isolators to validate geometry before committing to engineered units. \emph{Prototype only.} & Wire rope (proto) \\
\bottomrule
\end{tabularx}
\end{center}

\noindent\textbf{Primary picks for us:} \textbf{4$\times$ Enidine WR6-700-10} (or six, or step to WR8 for
harsh drops), mounted at $45^{\circ}$ with a rubber end-stop, because they damp the placement impact,
cover vibration and curb drops in the same part, need no air or power, and are maintenance-free. The
single open point is the \textbf{exact size/count}, which the supplier sizing tool fixes once the drop
height is known. Hold \textbf{Solution~F (guided deck + ACE SC$^2$)} as the upgrade for true drop loading.

% =====================================================================
\section{Priority of this modification}
\prio{\textbf{Priority: MEDIUM now, rising to HIGH before any payload-handling trials.}}

\begin{itemize}
\item \textbf{Not a bring-up blocker.} An empty or gently hand-loaded deck on a rigid mount is fine for
powering up the robot and validating electronics, navigation and the safety chain. It does not gate first
motion.
\item \textbf{Becomes critical for payload operations.} Once the robot does automated or repeated payload
placement --- and once it drives a real floor with joints, dock-plates and curb lips --- the rigid deck
funnels every placement and curb impact straight into the frame, fasteners and lift mechanism, where
repeated multi-kN spikes drive fatigue and shock the electronics. This must be resolved before field
trials.
\item \textbf{Ranking against other open items.} It sits \emph{below} the chassis-geometry and drivetrain
decisions that determine whether the platform can be built at all, but it \emph{interacts with the lift
mechanism} the deck mounts to --- so the suspension interface must be fixed together with the lift, not
bolted on afterwards.
\end{itemize}

\noindent\textbf{Recommended action:} issue an RFQ / order for 4--8$\times$ Enidine WR6-700-10 (confirm
exact part and count via the supplier sizing tool once drop height/velocity is measured), mount them at
$45^{\circ}$ with hard limit stops and a rubber end-stop, and keep the guided-deck + ACE SC$^2$ assembly
(Solution~F) as the fallback for genuine drop loading.

% =====================================================================
\begin{thebibliography}{99}
\small
\bibitem{abb} ABB. \textit{AMR P603V Product Specification} (3HKA00000226401) --- lift-type AMR; deck
raised on ball-screw lifters; positioning to $\approx\pm5$\,mm; on-deck load-measurement sensors.
\bibitem{enidine} Enidine (ITT). \textit{Wire Rope Isolators} --- stainless cable, RoHS aluminium bars;
qualified to MIL-STD-810, BV43-44, MIL-STD-167 and similar; used in automated-warehouse equipment.
\url{https://www.enidine.com/en-us/products/wire-rope-isolators}
\bibitem{uspto45a} USPTO. \textit{Multi-axis shock and vibration isolation system}, US~6,530,563
(wire rope isolators at $\approx 45^{\circ}$ to the suspended mass for balanced three-axis stiffness).
\url{https://image-ppubs.uspto.gov/dirsearch-public/print/downloadPdf/6530563}
\bibitem{uspto45b} USPTO. \textit{Shock resistant device}, US~7,260,476 (mounting wire rope isolators at
$\approx 45^{\circ}$ equalises compression/shear/roll stiffness and extends life).
\url{https://image-ppubs.uspto.gov/dirsearch-public/print/downloadPdf/7260476}
\bibitem{acesc} ACE Controls. \textit{SC$^2$ Self-Compensating Industrial Shock Absorbers} --- effective
weight 1--1{,}550\,kg, integrated positive stop, side-load adapter to $25^{\circ}$, react to changing
energy without adjustment. \url{https://www.acecontrols.com/us/}
\bibitem{aceheavy} ACE Controls / Valin. \textit{Heavy / self-compensating industrial shock absorbers}
(self-compensating units react to changing energy conditions without adjustment).
\url{https://www.valin.com/motion-control-and-automation/products/heavy-industrial-shock-absorbers-ace-controls-inc}
\bibitem{wr6} Enidine. \textit{WR6-700-10 wire rope isolator datasheet} (load capacity $\approx$50.8\,kg
/ 112\,lb; thickness $\approx$43\,mm).
\url{https://datasheets.globalspec.com/ds/itt-inc/wr6-700-10/36ff8e41-4c53-43df-81f6-df1d272cb7a7}
\bibitem{wr8} Enidine. \textit{WR8-400-08 wire rope isolator datasheet} (load capacity $\approx$68.9\,kg
/ 152\,lb; thickness $\approx$54\,mm).
\url{https://datasheets.globalspec.com/ds/itt-inc/wr8-400-08/07d5da4c-76cb-48fa-be8d-d096f6d5a05b}
\bibitem{lord} Parker LORD. \textit{Flex-Bolt Sandwich Mounts} (small group $\approx$31--645\,N
$\approx$3--66\,kg per mount; metric thread).
\url{https://www.modusadvanced.com/parker-lord-vibration-isolation-mounts}
\bibitem{firestone} Firestone Industrial. \textit{Airstroke / Airmount Engineering Manual \& Design Guide}
(air-spring isolators; highest isolation of any type; no spring chatter).
\url{https://www.firestoneairide.com/content/dam/fsip/pdfs/airstroke/Actuators-and-Isolators-Metric-Design-Guide.pdf}
\bibitem{acepal} ACE Controls. \textit{PAL-3 to PAL-9 low-frequency pneumatic levelling mounts}
(automatic level control; up to 385\,kg per air spring; natural frequency to $\sim$1.7\,Hz).
\url{https://www.acecontrols.com/us/products/vibration-control/low-freq-pneumatic-levelling-mounts.html}
\bibitem{mason} Mason Industries. \textit{Spring Isolators / housed spring mounts} (steel-coil low-frequency
vibration mounts; horizontal stiffness drops with deflection --- add limit stops).
\url{https://mason-ind.com/spring-isolators/}
\bibitem{casterhq} CasterHQ. \textit{Spring-loaded casters --- AGV guidance} (sprung compliance under a
deck introduces variable deck height that breaks sensor / dock-plate alignment unless travel is constrained).
\url{https://casterhq.com/collections/casters-spring-loaded}
\bibitem{pgfun} PGFUN. \textit{10-loop 304 stainless wire rope isolator} (low-cost prototype isolators;
more loops give better cushioning).
\url{https://www.amazon.com/PGFUN-Isolator-Stainless-Anti-Vibration-Absorber/dp/B09W695DJ6}
\end{thebibliography}

\end{document}
